\documentclass[11pt]{article}
\usepackage[a4paper,margin=27mm]{geometry}
\usepackage{amsmath,amssymb,amsthm,hyperref}
\newtheorem{theorem}{Theorem}\newtheorem{lemma}[theorem]{Lemma}
\newcommand{\F}{\mathbb F}\newcommand{\gen}[1]{\langle#1\rangle}
\newcommand{\PSL}{\operatorname{PSL}}\newcommand{\PSU}{\operatorname{PSU}}
\title{The two regular ends of a chiral marking}
\author{A recovery of Theorem 1.3 of Leemans--Vandenschrick}\date{}
\begin{document}\maketitle
We prove that neither $\PSL_3(q)$ nor $\PSU_3(q)$, for $q\ge2$,
is the automorphism group of a chiral polytope of rank greater than five.
The acquired classifications of regular subgroups give the two possible
local actions. The argument must allow one end to preserve a point
and the other a line. This distinction replaces an unjustified inference
in the source's even-characteristic proof.

We use the standard group characterization of chiral polytopes, in the
notation of Section~2 of the source, with $r\ge3$. Thus $G=\gen{\sigma_1,\ldots,
\sigma_{r-1}}$, every $\sigma_i$ is nonidentity, and
$\tau_{ij}=\sigma_i\cdots\sigma_j$ is an involution for $i<j$.
For $I\subseteq\{0,\ldots,r-1\}$ put
\[
 G^I=\gen{\tau_{ij}:i\le j,\ i-1,j\in I}.
\]
The actual intersection property is $G^I\cap G^J=G^{I\cap J}$;
in particular a singleton index set gives the trivial group. No
quotient intersection property will be used.

\section*{The regular ends and a proper normal subgroup}
\begin{lemma}\label{ends}
The involutions $\tau_{12},\ldots,\tau_{1,r-1}$ form a directly
regular string C-marking of rank $r-2$, of type
$[|\sigma_3|,\ldots,|\sigma_{r-1}|]$. There is a second such marking
$\tau_{r-2,r-1},\ldots,\tau_{1,r-1}$, with reversed type
$[|\sigma_{r-3}|,\ldots,|\sigma_1|]$.
\end{lemma}
\begin{proof}
For $J\subseteq\{2,\ldots,r-1\}$ the subgroup generated by
$\tau_{1j}$ for $j\in J$ is exactly $G^{\{0\}\cup J}$.
Indeed, each remaining generator with both endpoints in $J$ is
$\tau_{i,j}=\tau_{1,i-1}^{-1}\tau_{1,j}$. Actual IP therefore gives
all subset intersections of the proposed marking. For $k-j\ge2$,
$\tau_{1j}\tau_{1k}=\tau_{j+1,k}$ is an involution, proving distant
commutation. Adjacent products are $\sigma_{j+1}$, as asserted.

The even-word subgroup of any involutory marking is generated by its
adjacent products. Here it is $\gen{\sigma_3,\ldots,\sigma_{r-1}}$,
a normal subgroup of index at most two. It cannot contain $\tau_{12}$:
its intersection with $G^{\{0,2\}}$ is $G^{\{2\}}=1$.
Thus the index is two. Applying the same argument to the reversed
inverse rotation tuple gives the second marking. This also proves the
source's Lemma~2.3 directly from its actual intersection convention.
\end{proof}

\begin{lemma}\label{two}
If $r\ge3$ and some $\sigma_i$ has order two, then $G$ has a proper
nontrivial normal subgroup. Also $\sigma_i$ and $\sigma_j$ commute
whenever $|i-j|>2$.
\end{lemma}
\begin{proof}
First note an elementary consequence of the string relations. If
$B^2=(AB)^2=(BC)^2=(ABC)^2=1$, then
$BAB=A^{-1}$, $BCB=C^{-1}$, and
\[
 1=(ABC)^2=A(BCB)(BAB)C=AC^{-1}A^{-1}C.
\]
Taking $A=\sigma_i$, $B=\sigma_{i+1}\cdots\sigma_{j-1}$ and
$C=\sigma_j$ proves the distant commutation statement.

Now let $t=\sigma_i$ have order two, and put
$L=\gen{\sigma_1,\ldots,\sigma_{i-1}}$ and
$R=\gen{\sigma_{i+1},\ldots,\sigma_{r-1}}$.
Conjugation by $t$ inverts the adjacent generator on either side.
It centralizes every longer interval ending immediately before it or
starting immediately after it, because both that interval and its
product with $t$ are involutions. These intervals and the adjacent
generator generate the corresponding side, so $t$ normalizes $L,R$.
Distant commutation shows that the two sides commute except possibly
$\sigma_{i-1},\sigma_{i+1}$. The identity above, with $B=t$, covers
that pair as well. Moreover $L\cap R=1$ by their disjoint vertex-index
sets. Hence $H=LR$ is a normal direct product in $G$.

Here properness is an actual IP assertion, not an assumption of
independence of rotation generators. If $t=lr$ with $l\in L,r\in R$,
then
\[
 r\in\gen{L,t}\cap R
 =G^{\{0,\ldots,i\}}\cap G^{\{i,\ldots,r-1\}}=1.
\]
It follows that $t\in L\cap\gen{t}=G^{\{i-1\}}=1$, a
contradiction. The same calculation holds at an endpoint, with the
empty side trivial. Since $r\ge3$, at least one other nonidentity
rotation belongs to $H$, so $H$ is nontrivial.
\end{proof}

We acquire the complete even-characteristic regular subgroup theorem:
a nondegenerate regular string marking in $\PSL_3(F)$ has rank at
most four, and at rank four it has, up to coordinates and duality,
\[
 R_0=I+E_{23},\quad R_1=I+aE_{31},\quad
 R_2=I+bE_{13},\quad R_3=I+E_{21},\qquad ab\ne0.       \tag{1}
\]
We also acquire the complete odd-characteristic theorem. Its regular
rank bound is four without a nondegeneracy assumption. Its nondegenerate
directly regular rank-four markings are the Heisenberg marking
$[p,2p,p]$ and the simplex marking $[3,3,3]$ in characteristic five.
The other two conic markings have simple groups and are not directly
regular. These are actual marked-group statements, not just group orders.

The classical simplicity and order formulas for the projective linear
and unitary groups imply that $\PSL_3(q)$ is simple for $q\ge2$,
and $\PSU_3(q)$ is simple for $q\ge3$. Also $\PSU_3(q)$ embeds
in $\PSL_3(q^2)$: the intersection of the scalar centre of
$\operatorname{SL}_3(q^2)$ with $\operatorname{SU}_3(q)$ is precisely
the latter's centre. Lemmas~\ref{ends} and~\ref{two}, and the acquired
rank bounds, consequently give $r\le6$ for these simple groups.

\section*{The small unitary group}
\begin{lemma}\label{small}
In $\PSU_3(2)$ no two distinct involutions commute. In particular it
cannot be the group of a chiral polytope of rank at least five.
\end{lemma}
\begin{proof}
The scalar kernel has odd order three. Every projective involution
therefore has a unique involutory lift, and projectively commuting
involutory lifts commute: their scalar commutator has square one.
In characteristic two a nonidentity involution of dimension three is
$T=I+uf$, where $u\ne0$, $f\ne0$ and $f(u)=0$.
Let $h$ be the nonsingular Hermitian form, linear in its second argument.
For $x\in\ker f$, preservation of $h(x,y)$ by $T$ gives
$h(x,u)f(y)=0$ for every $y$. Thus $\ker f=u^\perp$, and
$f(x)=\lambda h(u,x)$ for a nonzero scalar $\lambda$.
The equation $f(u)=0$ makes $u$ isotropic; preservation of $h$ now
reduces to $\lambda+\lambda^2=0$. Over $\F_4$ it has just one
nonzero solution. Consequently each projective centre supports exactly
one nonidentity unitary elation.

Two distinct commuting elations in a projective plane have the same
centre or the same axis, by the acquired rank-one commutation lemma.
Here the axis is the Hermitian polar of the centre, so either condition
forces the same centre and hence the same involution. This proves the
first claim. A rank-three string C-marking would contain distinct
commuting first and third generators. Lemma~\ref{ends} therefore rules
out chiral rank at least five. We have not used, or re-proved, the source's
stronger historical computation of no chiral markings at any rank.
\end{proof}

\section*{The mixed point and line case}
Assume henceforth that a simple group above has a chiral marking of
rank six. Define the actual regular ends
\[
 R=(\tau_{12},\tau_{13},\tau_{14},\tau_{15}),\qquad
 S=(\tau_{45},\tau_{35},\tau_{25},\tau_{15}).             \tag{2}
\]
Both are nondegenerate by Lemma~\ref{two}. In even characteristic
normalize $R$ as in (1), using a single duality if necessary.
The first, second and last marks of $S$ are then
\[
 S_0=R_1R_3=I+(e_2+a e_3)e_1^t,\quad
 S_1=R_0R_3=I+e_2(e_1^t+e_3^t),\quad S_3=R_3.          \tag{3}
\]
The marks $S_0,S_3$ have the same axis and different centres.
The acquired rank-four normal form says that in its point-preserving
orientation the outer pair has the same centre; in the dual orientation
it has the same axis. Thus $S$ is in the latter orientation.
Its first and third marks share a centre, as do its second and fourth.
Writing $u=e_2+a e_3$, we obtain
\[
 S_2=I+u f,\qquad f\ne0,\qquad f(u)=0.                 \tag{4}
\]
These conclusions concern the actual axes and centres in (3); they do
not infer a common fixed point of the two ends.

Since $\sigma_1=R_3S_2$ and $\sigma_5=R_2R_3$, the displayed lifts are
\[
 \Sigma_1=I+E_{21}+u f,\qquad
 \Sigma_5=I+bE_{13}+E_{21}.                            \tag{5}
\]
Lemma~\ref{two} requires their projective commutation. The second lift
is unipotent. If $\Sigma_1\Sigma_5=\zeta\Sigma_5\Sigma_1$ for a
scalar $\zeta$, conjugacy of $\Sigma_5$ and $\zeta\Sigma_5$ forces
$\zeta=1$ by their eigenvalues. Thus the lifts themselves commute.
But the first row of $\Sigma_5\Sigma_1-\Sigma_1\Sigma_5$ is $abf$,
which is nonzero. This excludes even characteristic.

\section*{The two odd-characteristic possibilities}
The two ends in (2) have a common first type entry $|\sigma_3|$.
If one is a simplex marking, the ambient characteristic is five, and
its entry three differs from the Heisenberg entry five. Thus either
both are simplex markings or both are Heisenberg markings.

In the simplex case all five $\sigma_i$ have order three. The classical
rotation presentation of the Coxeter group $[3,3,3,3,3]$ identifies its
rotation group with $A_7$. Hence $G$ is a nontrivial quotient of the
simple group $A_7$ and has order $2520$. This is impossible here:
$|\PSL_3(2)|=168$, $|\PSL_3(3)|=5616$, and for $q\ge4$
\[
 |\PSL_3(q)|=\frac{q^3(q^3-1)(q^2-1)}{(3,q-1)}\ge20160.
\]
Similarly $|\PSU_3(2)|=72$, $|\PSU_3(3)|=6048$, and for $q\ge4$
its order is at least $4^3(4^3+1)(4^2-1)/3=20800$.

Finally suppose both ends are Heisenberg markings. Up to coordinates
and a single duality, the acquired normal form gives
\[\begin{aligned}
 R_0&=\begin{pmatrix}1&0&0\\0&-1&0\\0&0&-1\end{pmatrix},&
 R_1&=\begin{pmatrix}1&0&0\\1&-1&0\\0&0&-1\end{pmatrix},\\
 R_2&=\begin{pmatrix}-1&0&0\\0&-1&0\\0&1&1\end{pmatrix},&
 R_3&=\operatorname{diag}(-1,-1,1).
\end{aligned}\]
These are the unique involutory determinant-one lifts. Therefore
\[
 S_0=\begin{pmatrix}-1&0&0\\-1&1&0\\0&0&-1\end{pmatrix},\quad
 S_1=\operatorname{diag}(-1,1,-1),\quad S_3=R_3.
\]
The first two marks of $S$ have common centre $e_2$ and distinct axes.
In the displayed Heisenberg normal form the first pair has a common
axis and the last pair a common centre; in its dual these roles reverse.
Consequently the last pair of $S$ has common axis $x_3=0$.
This accounts for the dual orientation explicitly, and yields
\[
 S_2=\begin{pmatrix}-1&0&u\\0&-1&v\\0&0&1\end{pmatrix}.
\]
Commutation with $S_0$ gives $u=2v$, although the remaining contradiction
does not need this extra restriction. We have
\[
 \Sigma_1=R_3S_2=I-uE_{13}-vE_{23},\qquad
 \Sigma_5=R_2R_3=I-E_{32}.
\]
Again projective distant commutation lifts to ordinary commutation
because $\Sigma_5$ is unipotent. Its matrix difference is
\[
 \Sigma_1\Sigma_5-\Sigma_5\Sigma_1
       =uE_{12}+vE_{22}-vE_{33}.
\]
Thus $u=v=0$, making $\sigma_1$ trivial, a contradiction.
Together with Lemma~\ref{small} this completes the claimed rank bound.

\section*{Attribution and exact scope}
The nonexistence result is Theorem~1.3 of the fixed 2022 source, and
Lemmas~\ref{ends} and~\ref{two} reprove its Lemmas~2.3,6.1,6.2.
The actual regular subgroup classifications are separately acquired
proofs of Section~5. The even argument treats both point and line orientations; the odd
argument justifies the homology-axis inference and accounts for the
projective scalar in distant commutation. The source's simplex
phrase ``$q=5$'' is read as characteristic five for this subgroup
alternative. The rank bound includes $\PSU_3(2)$ by the elementary
unitary argument above, without importing an unspecified computer
search. The broader rank-three subgroup classifications, the source's
small negative chiral censuses, and its remaining unitary existence
questions are separate obligations, not consequences claimed here.
\end{document}
